\subsection{后端设计}

后端为 2 宽乱序执行流水：译码 $\rightarrow$ 重命名 $\rightarrow$ 分发 $\rightarrow$ 保留站发射 $\rightarrow$ 执行写回 $\rightarrow$ 提交。
与统一 Issue Queue 的顺序双发射不同，本设计使用\textbf{异构保留站}分离 ALU/MEM/MDU 发射策略，并以\textbf{ROB 编号}作为重命名标签，结果驻留 ROB，无独立物理寄存器堆。

\subsubsection{译码模块设计（ID）}

每拍从 IB 取出至多两条指令，例化两个组合译码器。
译码产生 FU 类型（ALU/MEM/MDU）、独热操作码、立即数、源与目的寄存器信息，以及 CSR、TLB、\texttt{cacop}、\texttt{ibar} 等特权向量；
同时检测 INE、SYS、BRK、IPE 等译码例外。
分支指令在译码阶段标记为 ALU 类，实际方向与目标在 ALU 中计算，最终于提交级与预测结果比对。

\subsubsection{重命名映射表（RAT）与重命名级（RA）}

RAT 为 32 项结构，记录各逻辑寄存器是否 busy 及其映射的 ROB 编号。
重命名级主要完成：

\begin{enumerate}
  \item 向 ROB 申请成对编号（奇偶双体，保证双发射整对整对分配）；
  \item 读 RAT/ARF 获得源操作数或生产者 ROB 索引；同拍双指令之间的 WAR/WAW/RAW 通过旁路修正；
  \item 若指令写寄存器，则更新 RAT 的 busy 与 ROB 映射；
  \item 将静态信息与重命名结果锁存，送往分发级。
\end{enumerate}

提交时，若退休指令仍是某逻辑寄存器的最新生产者，则清除对应 busy，完成释放。
误预测或异常冲刷时，RAT 随全局冲刷恢复到最后正确提交的那一拍的状态。

\subsubsection{分发与保留站发射（Dispatch / RS）}

分发单元从 ROB 读取尚未就绪的操作数，并将指令路由到三类保留站：

\begin{itemize}
  \item \textbf{RS\_ALU0 / RS\_ALU1}：各 4 项，\textbf{乱序}挑选就绪项，可同拍向两路 ALU 双发；
  \item \textbf{RS\_MEM}：4 项，FIFO 为主，并允许有限的 load 越过未就绪 load，每拍最多入队 1 条访存；
  \item \textbf{RS\_MDU}：2 项，严格 FIFO，服务乘除、CSR 读、\texttt{rdcnt}、\texttt{cpucfg}、\texttt{invtlb} 打包等。
\end{itemize}

源操作数就绪且目标 FU 可接收时发射。CSR 写、TLB 维护、异常等副作用不在发射级提交，而由提交级统一处理，实现精确异常。

\subsubsection{执行模块设计（EXE）}

\begin{itemize}
  \item \textbf{ALU0/ALU1}：单周期算术逻辑与分支判定，结果写回 ROB，并给出“是否跳转”标志位与跳转目标供提交比对；
  \item \textbf{LSU}：AGU + DCache 两级。AGU 计算虚地址并经 MMU 翻译、检查对齐/页例外；DC 级完成 Store Buffer 前递、Dcache 访问与 miss 处理。支持命中快速写回与 store$\rightarrow$load 安全前递；
  \item \textbf{MDU}：非流水多周期单元。乘法为 DSP 流水（约 3 拍），除法为 CLZ 快速除法（约 5--10 拍），CSR 读类指令在此取数，TLB 相关指令先在执行级占位，然后提交级真正维护。
\end{itemize}

四路 FU 写回 ROB 对应项并置 complete。访存推测 store 先进入 STQ（深度 16），按 ROB 序在提交后转入 Store Buffer。

\subsubsection{重排序缓存（ROB）}

ROB 共 32 项，奇偶双体环形队列，是重命名标签、结果缓存与提交序的核心。
各项保存 PC、执行结果、访存/特权信息、预测元数据与异常向量等。
dispatch 侧提供多读口补操作数，FU 侧多写口完成，commit 侧读取队头。
分配时清complete，写回时置位。冲刷时按提交点回收尾指针之后的投机项。

\subsubsection{提交模块设计（COMMIT）与冲刷}

提交级每拍最多退休两条指令，并在下列情况退化为单提交或停提交：后条为需串行的特权/\texttt{ll}/\texttt{sc}/屏障、前条异常或 \texttt{ertn}、仅一条就绪等。
退休时写 ARF、按需写 CSR/TLB、推入 Store Buffer、更新 LLBit，并向 BPU/FTQ 回送训练。

误预测、例外、\texttt{ertn} 以及 CSR/TLB/\texttt{cacop}/\texttt{ibar} 等导致的重取，均由提交级生成冲刷类型，经 \texttt{ctrl} 广播到前端与后端，实现\textbf{单一恢复点}。
这与在执行级立即重定向的设计不同，换取了精确异常与特权顺序的简洁性，并以更高的双提交率与预测精度弥补部分延迟。
